\section{Information-Guided Interventions}
\label{sec:improving-information}

We evaluate whether information changes improve the intended predictive outcome. An informative observation source can guarantee an information floor under the Fisher-reserve design \citep{mansouri_arai_2026_fisher_glasses} (Appendix~\ref{app:information-reserves}); training and acquisition require separate behavioral evaluation. For predictor parameters $w$, empirical task-risk objective $\widehat R_{\mathrm{base}}$ (ERM or group distributionally robust optimization, Group DRO), and penalty $\lambda_{\FG}$, training uses equally weighted groups:
\begin{equation}
\min_w\;\widehat R_{\mathrm{base}}(w)
+\lambda_{\FG}\operatorname{CVaR}_{\alpha,g}
\!\left[\bigl\{\lambda_{\min}(\widehat F_g^{\rid}(w))+\varepsilon\bigr\}^{-1}\right].
\label{eq:fg-training-compact}
\end{equation}
Coordinates follow Section~\ref{sec:computing-auditing}. Among the frozen-feature models whose validation worst-group accuracy is within two percentage points of the best, we select the coefficient with the largest validation certificate. This selection uses validation group labels.

\begin{table}[htbp]
\centering
\caption{Frozen-feature Waterbirds training. Means with nominal 95\% normal-approximation intervals ($\pm1.96$ standard errors) across five seeds; uncentered feature--Schur ridge certificates use $\alpha=0.95$. Validation selection follows Section~\ref{sec:improving-information}.}
\label{tab:frozen-training}
\small
\begin{tabular}{@{}lccc@{}}
\toprule
Objective & \shortstack{Average\\accuracy (\%)} & \shortstack{Worst-group\\accuracy (\%)} & \shortstack{Fisher-Glass\\certificate} \\
\midrule
ERM & $77.2\pm0.1$ & $55.7\pm0.3$ & $0.198\pm0.001$ \\
ERM + Fisher-Glass & $80.4\pm0.2$ & $56.3\pm0.3$ & $0.399\pm0.006$ \\
Group DRO & $80.9\pm0.3$ & $67.3\pm0.3$ & $0.233\pm0.003$ \\
Group DRO + Fisher-Glass & $80.6\pm0.9$ & $69.4\pm1.3$ & $0.614\pm0.059$ \\
\bottomrule
\end{tabular}
\end{table}

Fisher-Glass approximately doubles ERM's certificate; with Group DRO it raises both the certificate and mean worst-group accuracy (Table~\ref{tab:frozen-training}). Stronger penalties can reduce accuracy (Appendix~\ref{app:extended-training}). These source-performance gains do not establish a repair of background sensitivity.

\FloatBarrier

In prospective fine-tuning, the validation certificate increases from $0.289$ to $0.330$, while controlled swap drops remain $52.9$ points for FG and $52.6$ for ERM. A separate Fisher-based frozen-slice repair under validation-accuracy constraints selects the unchanged model in all five seeds. On another cohort, post hoc selection raises rediscovered-slice information but reduces worst-group accuracy by $12.6$ points; selected coefficients match validation-NLL minimization in every seed (Appendix~\ref{app:extended-training}). Thus information improvement and the intended behavioral improvement are separate outcomes.

For candidate $a$ with modified feature/input $x_a^{\mathrm{cf}}$ and label $y_a$, \emph{tail-transverse influence} $\TTI(a)=G_\alpha(a)\loss_w(x_a^{\mathrm{cf}},y_a)$ weights current predictive loss by the first-order inverse-information CVaR decrease $G_\alpha(a)$ along $a$'s allocation path. This differentiates the Schur complement along its weakest task direction with tail-sensitive weights (Appendix~\ref{app:tail-transverse-influence}). Waterbirds ranks candidates by $\widetilde{\TTI}_\alpha(a)=\widetilde G_\alpha(a)\loss_w(x_a^{\mathrm{cf}},y_a)$, using the tail-weighted squared-residual surrogate $\widetilde G_\alpha$ defined in Appendix~\ref{app:feature-acquisition}.

Waterbirds candidates translate features by the within-class difference of background-group means, scaled by $\gamma_{\mathrm{cf}}$. Training data fit means and the nuisance probe; validation supplies candidates for retraining, and evaluation is held out. Budgets are fractions of original training size. At $\gamma_{\mathrm{cf}}=1$ and a $1\%$ budget, $\widetilde{\TTI}_\alpha$ yields a mean worst-group accuracy gain of $4.52$ points versus $3.27$ for high-loss selection; high loss produces the larger mean certificate gain. At this strength, $\widetilde{\TTI}_\alpha$ has the larger mean worst-group accuracy gain at budgets of $1\%$, $2\%$, and $5\%$; selection by $\widetilde G_\alpha$ alone yields smaller improvements (Appendix~\ref{app:feature-acquisition}).
